% Two independent derivatives and exact polynomial quadrature comparison.
\clearpage
\section*{Verification case: three meanings of numerical change}
\noindent A symbolic derivative, automatic differentiation of the
implemented operations and a finite-difference approximation do not
have identical error mechanisms. Consider
\[
F(x,y)=x^2y+\sin x.
\]
Its analytic gradient is
\(\nabla F=(2xy+\cos x,\;x^2)\).
At \((x,y)=(1,2)\), this is
\((4+\cos1,\;1)\), and its directional derivative along
\(d=(0.1,-0.2)\) is
\[
dF[d]=0.1(4+\cos1)-0.2\approx0.25403023.
\]
Forward-mode AD should reproduce the derivative of the implemented
smooth computation to ordinary rounding accuracy. A central
finite difference,
\[
\frac{F(1+hd_x,2+hd_y)-F(1-hd_x,2-hd_y)}{2h},
\]
has truncation error of order \(h^2\) for a sufficiently smooth
function, but its \emph{floating-point} error eventually grows as
\(h\) becomes too small. The step is a numerical design parameter,
not a universal magic constant.

Integration supplies a different error source. For \(f(x)=x^2\)
on \([0,1]\), the exact integral is \(1/3\). A single-interval
trapezoidal rule gives \(1/2\), a midpoint rule gives \(1/4\),
and Simpson's rule gives the exact \(1/3\) in exact arithmetic
because it integrates cubic-or-lower polynomials exactly.
For nonsmooth integrands or singular endpoints, do not apply these
polynomial error expectations without checking assumptions.

\paragraph{Reconstruction.}
Implement dual numbers for a scalar expression, two finite-difference
schemes, and three quadrature rules. Sweep the finite-difference step
on a log axis, record the error against the analytic directional
derivative, then test quadrature convergence on both a polynomial
and a non-polynomial function. Report function evaluation counts
alongside error, and distinguish a mathematically exact polynomial
rule from the rounded result computed in software.
